\documentclass[journal]{IEEEtran}

\usepackage{amsmath,amssymb,amsfonts}
\usepackage{authblk}
\usepackage{graphicx}
\usepackage{booktabs}
\usepackage{array}
\usepackage{microtype}
\usepackage{xcolor}
\usepackage[ruled,linesnumbered]{algorithm2e}
\usepackage[hidelinks]{hyperref}

\graphicspath{{figures/}}

\newcommand{\revps}{\ensuremath{\,\mathrm{rev/s}}}
\newcommand{\fsenv}{\ensuremath{f_{\mathrm{e}}}}
\newcommand{\snr}{\ensuremath{\mathrm{SNR}}}
\newcommand{\Jobj}{\ensuremath{J}}

% ---------------------------------------------------------------------------
% DRAFT MARKERS.
% \pending{...} flags a number or claim whose measurement has not yet landed.
% Every such site states what will be measured.
% \wip{...} flags a passage of text which is structurally in place but whose
% content depends on pending experiments and will be rewritten around the
% measured numbers.
% Remove both macros (and every use) before submission.
% ---------------------------------------------------------------------------
\newcommand{\pending}[1]{{\color{gray}\textsf{[pending:~#1]}}}
\newcommand{\wip}[1]{{\color{gray!70!black}\textit{[WIP:~#1]}}}

% Placeholder for a figure whose graphic is not yet generated.
\newcommand{\figplaceholder}[2]{%
  \framebox[\columnwidth]{\rule{0pt}{#1}\parbox{0.9\columnwidth}{\centering
      \pending{#2}}}}

\begin{document}

\title{Joint Decomposition of Quadcopter Ego-Noise into Harmonic and
  Broadband Components, with a Built-In Measure of Fit}

\author[1]{Dmitrii~Mukhutdinov \thanks{d.mukhutdinov@qmul.ac.uk}}
\author[1]{Lin~Wang \thanks{lin.wang@qmul.ac.uk}}

\affil[1]{School of Electronic Engineering and Computer Science, Queen
    Mary University of London}

\maketitle

\begin{abstract}
Noise emitted by multirotor unmanned aerial vehicles (drones) has a
distinct structure: for each rotor, both the aerodynamic noise of the
propeller and the mechanical noise of the motor are predominantly tonal,
with fundamental frequencies proportional to the rotor's time-varying
rotational speed, on top of a colored broadband component produced by
turbulence. Separating a recording into these two parts --- given the
rotor speed trajectories --- and, conversely, recovering the trajectories
from audio alone, are two faces of one problem, and both are useful: the
decomposition provides clean per-harmonic amplitude targets for
data-driven noise models, while audio-only speed recovery enables
post-annotation of the large body of drone audio for which no telemetry
exists. Given how clearly structured the noise is, the problem seems
easy; in practice, on real free-flight recordings, it is not. Four rotors
with close and intersecting speed trajectories, per-harmonic phase noise
that destroys coherence within fractions of a second at high harmonic
orders, and a strong broadband component defeat both classical
least-squares harmonic decomposition and existing tacholess order
tracking methods --- and, as we demonstrate experimentally, they defeat
them in a treacherous way: a plain least-squares residual actively
prefers physically wrong trajectory sets to the true ones. We propose a
joint probabilistic model of the harmonic and broadband components in
which phase noise enters through the increments of instantaneous
frequency and the broadband floor is an explicit smooth spectral surface
that charges a likelihood rent for the energy it absorbs. Maximum a
posteriori estimation in this model is performed by a block-coordinate
solver, and the one machine yields three readouts: the decomposition
itself, a refinement of the input speed trajectories, and a principled
scalar measure of how well a trajectory set fits the audio. We use the
measure to drive trajectory search, refining imprecise telemetry and
recovering speeds blind, and evaluate the whole system on indoor and
outdoor free-flight data against telemetry, an independent tachometer
track, and a downstream neural noise-generation task.
\end{abstract}

\begin{IEEEkeywords}
Ego-noise, drone audition, harmonic decomposition, order tracking,
Vold--Kalman filter, phase noise, rotor speed estimation, maximum a
posteriori estimation.
\end{IEEEkeywords}

\section{Introduction}
\label{sec:intro}

\IEEEPARstart{A}{coustic} sensing on board of a flying multirotor drone
is dominated by the drone's own \emph{ego-noise}. Any downstream task ---
speech enhancement, sound source localization, acoustic monitoring ---
benefits from a model of this noise, and the most informative side
channel for such a model is the set of rotor speed trajectories: the
rotational speeds determine the fundamental frequencies of every tonal
component in the noise and reflect the flight dynamics that shape the
broadband component. Recent data-driven noise models and noise-aware
enhancement systems consume rotor speeds directly as conditioning
input~\cite{yen2023ncm,gulli2025rotorcond}.

This creates a data problem. Neural models want large training corpora,
but free-flight drone audio annotated with per-rotor speed telemetry is
scarce: to our knowledge, the only public dataset is
DREGON~\cite{strauss2018dregon}, and our own outdoor recordings add a few
more flights. Meanwhile, unannotated drone audio is plentiful. If rotor
speeds could be recovered from audio alone --- precisely enough to serve
as training targets and conditioning inputs --- every unannotated
recording would become training data. The prospect looks reasonable to
the naked eye: on a spectrogram of a drone recording, the harmonic combs
of the rotors are often clearly visible, and one can trace their motion
by hand. It is tempting to conclude that the extraction task is easy.

It is not, and the reasons are instructive. First, a quadcopter carries
four rotors whose speeds in cruise are distinct but close --- on DREGON,
twin pairs sit $0.4$--$0.8\revps$ apart --- and their harmonic combs
interleave and cross. Second, the tonal components are not clean
sinusoids: each harmonic carries phase noise, and the coherence time
shrinks rapidly with harmonic order, so that above roughly the tenth
harmonic the ``line'' on the spectrogram is really a band-limited random
process that merely stays centered on the comb. Third, the broadband
component is strong, colored, and time-varying, and at low frequencies it
does not even add up power-wise across rotors in a way that would let one
attribute it to individual sources. Existing tacholess order-tracking and
multi-pitch methods, built for single shafts or for well-separated
harmonic sources, do not survive this combination
(Sec.~\ref{sec:related}).

There is also a fourth reason, which we only understood after building
and instrumenting a full least-squares decomposition pipeline, and which
we believe is the central methodological lesson of this paper. The
natural formalization of ``how well does a trajectory set fit the
audio?'' --- project the signal onto the harmonic subspace spanned by the
hypothesized combs and measure the residual --- is not merely imprecise;
it is \emph{actively misleading}. A projection residual rewards spectral
coverage, and coverage is free: every band the hypothesis opens can only
absorb more energy. On real recordings whose combs are weak and
decohered at high orders, we measured trajectory sets that are wrong by
$6$--$7\revps$ --- physically meaningless ``fans'' of closely spaced
tracks that tile the spectrum --- \emph{winning} the least-squares
objective against the true trajectories by a wide margin
(Sec.~\ref{sec:whynotls}). A search driven by such an objective converges
confidently to nonsense. Before one can search for rotor speeds, one must
first possess a measure of fit that ranks the truth first; and, as we
show, obtaining such a measure requires modeling the noise honestly ---
the harmonic part with its phase-noise structure, and the broadband part
as an explicit component that pays for the energy it explains.

We therefore pose the problem jointly. Our research questions are:
\begin{enumerate}
\item Given rotor speed trajectories, can we decompose a real recording
  into its harmonic and broadband components, such that the broadband
  residual is genuinely free of harmonic structure?
\item Can we construct a scalar measure of fit between a trajectory set
  and the audio which is reliable --- in particular, which ranks true
  trajectories above physically wrong ones where the least-squares
  residual provably fails?
\item Using such a measure, can we iteratively approach the true
  trajectories --- both from imprecise telemetry and from no
  initialization at all?
\end{enumerate}
The thesis of this paper is that a single model answers all three at
once. Maximum a posteriori (MAP) estimation under a joint harmonic-plus-broadband
noise model is performed by one block-coordinate solver, and the three
answers are three readouts of the same machine: the decomposition is its
converged latent state, the trajectory refinement is one of its
coordinate blocks, and the measure of fit is its converged objective
value.

Our contributions are as follows:
\begin{enumerate}
\item We formulate a probabilistic model of multirotor ego-noise in
  which per-harmonic phase noise enters through the \emph{increments} of
  instantaneous frequency --- so that harmonic lines decohere in phase
  yet remain anchored to the comb in frequency, as observed --- and the
  broadband component is a smooth time-varying spectral floor. We
  justify the model's structure (which noise terms are shared across
  harmonics and rotors, and which are individual) by direct phase
  coherence measurements on real recordings
  (Sec.~\ref{sec:model}).
\item We derive the MAP estimation problem for this model and show that
  its objective, unlike a projection residual, charges a ``rent'' for
  spectral coverage; we demonstrate on real data that this difference is
  not academic, by exhibiting the failure mode in which least squares
  prefers wrong trajectories and \pending{J-rescoring experiment}
  verifying that the proposed objective ranks them correctly
  (Secs.~\ref{sec:map}, \ref{sec:res-measure}).
\item We describe the block-coordinate solver and its three readouts.
  Its trajectory block is an information-form Kalman smoother over
  phase increments --- a principled descendant of Vold--Kalman order
  tracking that we relate, via an explicit capture-range analysis, to
  the classical VK frequency loop and its published adaptive variants
  (Secs.~\ref{sec:solver}, \ref{sec:capture}).
\item We evaluate on indoor (DREGON) and outdoor (Matrice~100) free
  flight along a single axis of initialization quality --- from oracle
  telemetry through corrupted telemetry to fully blind --- with an
  auditing methodology that treats the reference labels themselves as
  measurements: a telemetry clock-and-scale audit, an independent
  tachometer track, and a downstream neural noise-generation task that
  measures whether refined trajectories are \emph{better training
  labels} than the telemetry they started from
  (Secs.~\ref{sec:setup}, \ref{sec:results}).
\end{enumerate}

The organization of the paper is as follows. Section~\ref{sec:related}
reviews related work. Section~\ref{sec:model} builds the signal and noise
model and presents the coherence measurements that fix its structure.
Section~\ref{sec:map} derives the MAP problem, discusses why the
least-squares alternative fails, and describes the block-coordinate
solver. Section~\ref{sec:search} describes the global search stage used
for blind operation and the capture-range analysis that dictates the
architecture. Sections~\ref{sec:setup} and~\ref{sec:results} present the
experimental setup and results, Section~\ref{sec:discussion} discusses
limitations, and Section~\ref{sec:conclusion} concludes.

\section{Related Work}
\label{sec:related}

\emph{VK order tracking and adaptive variants.} The Vold--Kalman (VK)
filter was introduced for high-slew order
extraction~\cite{vold1993high}, extended to multiple independent
shafts~\cite{vold1997multi}, and characterized
in~\cite{herlufsen1999characteristics,tuma2000characteristics}; adaptive
reformulations followed~\cite{pan2007adaptive,pan2010extended}. Two
recent members split on the question that our capture-range analysis
turns on: where the instantaneous frequency (IF) is estimated. Li
\emph{et al.}~\cite{li2020vkswt} estimate IF \emph{outside} the filter,
by synchrosqueezed-wavelet ridge detection, and hand it to VK as a fixed
parameter; a ridge detector has essentially unbounded capture range, so
this architecture never meets the failure we analyze.
IAVKF~\cite{jiang2024iavkf} estimates IF \emph{inside} the loop: the
envelope-phase derivative is smoothed by a second, IF-domain VK problem
and added to the carrier, while a per-component bandwidth adapts to an
orthogonality fixed point. IAVKF peels components greedily and states
that it cannot handle overlapping IFs --- precisely the quadrotor twin
pair --- and its wide-to-narrow bandwidth adaptation carries no
dependence on harmonic order. Section~\ref{sec:capture} reads that
adaptation as an implicit capture schedule and makes it explicit.

\emph{Ridge extraction and tacholess order tracking.} Adaptive ridge
extraction~\cite{iatsenko2016ridges,aare2023ridge} and ridge-path
regrouping at crossings~\cite{chen2017rprg} are the blind front-end
competitors; tacholess angular-resampling methods recover the shaft rate
from the signal itself~\cite{peeters2019tacholess,wang2019tlot}, with
generalized demodulation~\cite{olhede2005generalized} and the fan-chirp
transform~\cite{weruaga2007fanchirp} as transform-domain relatives.
These methods target a single shaft; a single-component ridge detector
cannot prefer a shaft comb over its blade-passage octave, and it cannot
keep four near-coincident combs apart. This is why our global stage
scores whole comb hypotheses rather than individual lines.

\emph{Harmonic-plus-noise signal models.} The decomposition of a signal
into a deterministic harmonic part and a stochastic residual has a long
history in audio analysis and synthesis: spectral modeling synthesis and
the harmonic-plus-noise model
\pending{cite: Serra \& Smith SMS; Stylianou HNM}, and, recently,
differentiable digital signal processing, where a neural network
predicts per-harmonic amplitudes and a filtered-noise floor end to end
\pending{cite: Engel et al.\ DDSP}. On the estimation side, Bayesian and
maximum-likelihood models of harmonic signals in colored noise are well
developed for pitch estimation
\pending{cite: Christensen \& Jakobsson; Nielsen et al.\ fast NLS pitch},
and smooth spectral floors are classically estimated by
Whittle-likelihood and cepstral methods
\pending{cite: Whittle; cepstral PSD smoothing}. Our model belongs to
this family, with three features the application forces on it:
per-harmonic phase noise expressed in IF increments (rather than
stationary sinusoids or slowly-varying deterministic phase), joint
treatment of multiple closely-spaced combs with shared shaft-level
errors, and the explicit use of the MAP objective as a \emph{measure}
that drives trajectory search. We are not aware of prior work that uses
the failure of the projection residual as a motivation for the joint
model; the closest in spirit are penalized and marginal-likelihood
arguments for model-order selection \pending{cite: model selection /
Occam factor reference}.

\emph{Drone audition.} Ego-noise suppression spans
BSS~\cite{wang2020bss}, learned masking~\cite{wang2020dnntf}, and deep
enhancement; several systems consume rotor
speeds~\cite{yen2023ncm,gulli2025rotorcond}. Our own earlier work
approached the decomposition by frame-wise least-squares projection onto
demodulated harmonic subspaces; the failure analysis of that approach
(Sec.~\ref{sec:whynotls}) is what motivated the present model. To our
knowledge no prior work recovers all four rotor rates blind at twin
separations below $1\revps$, and no prior work on drone ego-noise
provides a validated decomposition of real free-flight recordings into
harmonic and broadband parts.

\section{Signal and Noise Model}
\label{sec:model}

In this section we build the generative model of the recorded signal.
We proceed in three steps: the deterministic skeleton
(Sec.~\ref{sec:model-skeleton}), the phase noise structure
(Sec.~\ref{sec:model-phasenoise}), and the broadband floor
(Sec.~\ref{sec:model-floor}). Each modeling decision that is not forced
by first principles is fixed by a measurement on real data; the
measurements are summarized in Sec.~\ref{sec:model-measurements}.

\subsection{The deterministic skeleton}
\label{sec:model-skeleton}

Audio $x_c[n]$ ($C \le 8$ microphones, sample rate $f_s = 16$\,kHz,
sample spacing $\Delta = 1/f_s$) is recorded on a quadcopter whose rotor
$i \in \{1,\dots,4\}$ spins at $r_i(t)$ revolutions per second
($74$--$91\revps$ in cruise on our data). Both the mechanical noise of
motor $i$ and the tonal aerodynamic noise of its propeller are harmonic
with fundamental frequency $r_i(t)$; a propeller with $B$ blades
concentrates aerodynamic energy on the multiples of the blade-passage
frequency $B\,r_i(t)$, which are a subset of the shaft harmonics, so a
single comb per rotor with per-harmonic amplitudes covers both. Given a
rate estimate $\hat r_i$ (from telemetry or from a search stage), define
the \emph{carrier phase} as the cumulative sum
\begin{equation}
  \hat\varphi_i[n] = 2\pi \Delta \sum_{m \le n} \hat r_i[m].
  \label{eq:carrier}
\end{equation}
Every phase in this paper is an integral of a rate; this is not a
notational preference but the load-bearing structural fact of the model,
as the next subsection explains. The observation model is
\begin{equation}
  x_c[n] = \sum_{i,k}
  \operatorname{Re}\Big[ A_{c,i,k}[n]\,
    e^{\,\mathrm{j} ( k \hat\varphi_i[n] + k\theta_i[n] +
      \psi_{i,k}[n] )} \Big] + n_c[n],
  \label{eq:model}
\end{equation}
where $A_{c,i,k}[n]$ are slowly varying complex envelopes (their
argument absorbs the unknown initial phase of each harmonic at each
microphone), $\theta_i$ and $\psi_{i,k}$ are the phase-noise terms
defined next, and $n_c$ is the broadband component. We use $K = 80$
harmonics, the largest number that keeps every line below the Nyquist
frequency at our rotor speeds.

In words: the model says that once the carrier is right, what remains of
each harmonic is a narrowband complex envelope; everything that is not
narrowband around some comb line belongs to the floor.

\subsection{Phase noise lives in rate increments}
\label{sec:model-phasenoise}

Two facts about real recordings must be reconciled. First, the harmonic
lines \emph{decohere}: the phase of harmonic $k$, tracked against a
perfect carrier, drifts without bound, and the drift is faster for
larger $k$ --- above $k \approx 10$ on our data the coherence time drops
below one second, so no method that integrates absolute phase coherently
over long spans can work there. Second, the lines \emph{do not wander
off the comb}: the instantaneous frequency of harmonic $k$ stays
anchored at $k\,r_i(t)$, visibly, for the whole recording. A model of
phase noise as a random walk in \emph{phase} captures the first fact but
violates the second (a phase random walk is a white-noise IF, which a
physical rotor with inertia cannot produce and our telemetry does not
show); a model of phase noise as small pointwise phase deviations
captures neither. The reconciliation is to put the noise in the
\emph{rate} domain and integrate:
\begin{equation}
  \theta_i[n] = 2\pi\Delta \sum_{m\le n} \delta_i[m],
  \qquad
  \psi_{i,k}[n] = 2\pi\Delta \sum_{m\le n} \nu_{i,k}[m],
  \label{eq:phase-integrals}
\end{equation}
where $\delta_i$ is the \emph{shaft rate error} of rotor $i$ (the true
rate minus the carrier's rate; it multiplies $k$ in \eqref{eq:model}
because a shaft speed deviation shifts harmonic $k$ by $k$ times as
much) and $\nu_{i,k}$ is a small per-harmonic rate wobble that models
everything the rigid comb misses (blade-to-blade asymmetries,
aeroacoustic modulation). Both are modeled as stationary, bounded,
temporally correlated processes --- concretely, Ornstein--Uhlenbeck
processes with variance $\sigma^2$ and correlation time $\tau_c$. The
integral of a bounded stationary rate is a diffusing phase: over spans
short against $\tau_c$ the phase performs a random walk with diffusion
rate proportional to $\sigma^2 \tau_c$, which produces exactly the
observed pair of facts --- unbounded phase drift (decoherence, with
Lorentzian-type line broadening growing with $k$) together with bounded
frequency excursion (lines anchored to the comb).

Two structural questions remain: which of these terms are shared, and
how strong is each. Both are answered by measurement
(Sec.~\ref{sec:model-measurements}); here we state the outcome. The
shaft term $\delta_i$ is \emph{rig-common} to a large degree --- the
low-frequency drift of all four rotors against telemetry moves together
--- which motivates splitting it into a common component and small
per-rotor components in the solver. The per-harmonic wobble $\nu_{i,k}$
grows with $k$: the effective linewidth of harmonic $k$ scales
approximately linearly in $k$, at roughly $0.6k$\,Hz on our recordings,
which partitions the comb into three regimes: low harmonics that stay
phase-coherent over seconds, a middle band that is coherent only over
fractions of a second, and high harmonics that are best described as
narrowband noise processes centered on the comb. A propagation-delay
term (per rotor-microphone pair, scaling the phase by $k$ like the shaft
term) belongs in \eqref{eq:model} on physical grounds; we carried it
through the derivation, measured it to be negligible against the other
terms on our arrays, and pin it to zero in all experiments.

\subsection{The broadband floor}
\label{sec:model-floor}

The broadband component $n_c$ is colored, with a smooth and smoothly
time-varying spectral shape. We model it as a Gaussian process, independent
across microphones, with log power spectral density $\log S_c(f, t)$
constrained to be smooth in both arguments. Two remarks fix the scope of
this choice. First, smoothness is the load-bearing property: the floor
must be flexible enough to follow the true broadband timbre but rigid
enough that it \emph{cannot} imitate a comb --- the whole discipline of
the decomposition rests on the floor being unable to absorb harmonic
structure and the comb being unable to absorb the floor
(Sec.~\ref{sec:map}). Second, we deliberately do not attribute broadband
energy to individual rotors. We attempted this attribution using
multichannel data and known geometry and found it unidentifiable in
steady flight: on DREGON the microphones sit in the rotor wash and
inter-microphone coherence of the broadband component is near zero
(median $0.014$ against $0.98$ for an ideal diffuse field at low
frequency), the four rotor speeds are nearly collinear in cruise
(variance inflation factors $21$--$118$), and measured per-rotor powers
do not add up below $1$\,kHz (joint operation exceeds the sum of
solo-rotor operation by up to $10.9$\,dB). The model therefore carries
one floor per microphone, with a single common flight-state driver, and
makes no per-rotor broadband claim.

\subsection{What the measurements say}
\label{sec:model-measurements}

\wip{This subsection summarizes the coherence-probe protocol and its
outcomes; the figures will be regenerated for the paper from the
existing probe results.} Three measurements on real recordings fix the
model structure stated above.

\emph{Decoherence versus $k$.} \pending{figure: per-$k$ coherence time /
linewidth measurement on DREGON and Matrice~100, showing the
approximately linear-in-$k$ linewidth $\approx 0.6k$\,Hz and the three
coherence regimes.}

\emph{Rig-commonness of the shaft drift.} Low-pass filtered,
$k$-normalized phase increments of strong harmonics are correlated
across rotors nearly as strongly as across harmonics of the same rotor
(cross-rotor correlation $0.46$ against same-rotor $0.47$ at $2$\,s
smoothing on Matrice~100), identifying the dominant slow drift as
rig-common rather than per-rotor. \pending{figure: pair-class
correlation versus smoothing time constant.}

\emph{Arrival term.} The per-(rotor, microphone) propagation delay
predicted by array geometry is measurable on the outdoor rig
(cross-channel comb-phase delay slope $1.013$ against a predicted $1.0$,
window-to-window repeatability $0.098$\,ms) and contributes $-0.92$\,ms
to the effective carrier alignment --- two orders of magnitude below the
phase-noise terms at the harmonics that matter, which justifies pinning
it to zero.

\section{Decomposition as MAP Estimation}
\label{sec:map}

We now turn the model into an estimator. Section~\ref{sec:map-objective}
states the MAP problem and reads its terms;
Sec.~\ref{sec:whynotls} explains, and demonstrates on real data, why the
seemingly simpler least-squares alternative must be rejected;
Sec.~\ref{sec:solver} describes the block-coordinate solver;
Sec.~\ref{sec:readouts} states the three readouts and their relation to
classical order tracking.

\subsection{The objective}
\label{sec:map-objective}

Collect the unknowns: the complex envelopes $A = \{A_{c,i,k}\}$, the
phase-noise trajectories $\theta = \{\theta_i\}$ and
$\psi = \{\psi_{i,k}\}$, and the floor $S = \{S_c(f,t)\}$. Within an
analysis window (12\,s in our implementation) the correlation time of
the rate noise exceeds the window, so the OU priors reduce to their
local limit: white noise on the \emph{rate increments}, i.e.\ a
second-difference (curvature) penalty on the integrated phases. Writing
$D_2$ for the second-difference operator on the envelope-rate time grid
and $P_{c}(f,t)$ for the short-time residual power of microphone $c$
after subtracting the harmonic reconstruction of \eqref{eq:model}, the
negative log posterior is, up to constants,
\begin{multline}
  \Jobj(A, \theta, \psi, S) =
  \sum_{c,f,t} \left[ \frac{P_c(f,t)}{S_c(f,t)} + \log S_c(f,t) \right]
  \\
  + \lambda_\theta \sum_i \| D_2 \theta_i \|^2
  + \sum_{i,k} \lambda_\psi(k)\, \| D_2 \psi_{i,k} \|^2
  \\
  + \sum_{c,i,k} \rho_k^2\, \| D_2 A_{c,i,k} \|^2
  + R(\log S),
  \label{eq:J}
\end{multline}
where $R(\log S)$ is the smoothness penalty on the log-floor and the
penalty weights map one-for-one to bandwidths by the standard VK
relation: a target $-3$\,dB bandwidth $b$ at envelope rate $\fsenv$
corresponds to $\lambda(b) = \big(2 \sin(\pi b / \fsenv)\big)^{-4}$.
This single calibration relation is shared by every smoothness term in
the system --- envelopes, shaft phases, per-track phases --- so the
hyperparameters of \eqref{eq:J} are bandwidths, all in hertz, all
interpretable, and all frozen across recordings
(Table~\ref{tab:constants}).

Read the first sum carefully, because it is where this objective parts
ways with least squares. For fixed floor $S$, the data term is a
\emph{whitened} residual: energy is counted in units of the local floor.
Residual energy at floor level contributes $P/S \approx 1$ per cell
regardless of what the comb does, so a harmonic track gains nothing by
absorbing floor-level energy --- the floor already explains it, at the
price of the $\log S$ term, which is the ``rent'': a floor that inflates
to absorb everything pays $\log S$ for every cell, and the optimal
stationary point balances the two ($S^\ast = P$ cell-wise in the
unpenalized limit, which is the maximum-likelihood variance estimate).
A harmonic track earns likelihood only where it explains energy
\emph{above} the local floor --- that is, on actual lines --- and it
pays its own curvature rent through $\rho_k^2$. The smoothness penalty
$R$ prevents the converse theft: a spiky floor could otherwise drop onto
each comb line and claim it; a smooth floor cannot.

\subsection{Why not least squares}
\label{sec:whynotls}

The alternative that every practitioner tries first --- we did, at
length --- is to skip the floor and the phase-noise structure: fix the
trajectories, project the signal onto the harmonic subspace (a bank of
demodulated narrowband envelopes, solved jointly as penalized least
squares), and use the residual energy as the measure of fit. This is the
profiled least-squares residual; the second-generation VK filter is
exactly its per-window solver. It is attractive because it is convex in
the envelopes, exactly solvable, and fast. It fails, and the failure is
structural: the projection residual rewards \emph{coverage}. Every
narrowband subspace the hypothesis opens can only absorb energy; nothing
charges for opening bands in the wrong places. When the comb is strong,
the true hypothesis wins anyway, because bands on real lines absorb far
more than bands on the floor. When the comb is weak and decohered at
high orders --- the regime of every real drone recording past
$k \approx 10$ --- coverage beats correctness.

We measured this failure directly. On three DREGON cruise windows we
compared the profiled least-squares fitness of the telemetry trajectories
against trajectory sets produced by unconstrained multi-start
optimization of that same fitness. On all three windows the optimizer
found solutions that \emph{win the objective} by a wide margin
($0.609$--$0.672$ against telemetry's $0.695$--$0.727$; lower is better)
while sitting $6.5$--$7.2\revps$ from the telemetry in
permutation-matched MAE. The winning configurations are physically
meaningless: single tracks that swallow a whole twin pair, plus ``fans''
of smooth, closely spaced extra tracks whose interleaved harmonics tile
the upper spectrum. On the outdoor recordings, whose combs are strong
and well resolved, the same objective ranks telemetry first and the
optimizer merely fails to find it --- confirming that the defect is not
in the optimizer but in what the objective rewards on weak-comb data.
\pending{figure: one DREGON window --- spectrogram, telemetry comb
versus fan-solution comb overlay, and both trajectory sets; table:
fitness and PIT-MAE of telemetry / fan / blind chain on the three
windows (numbers measured, to be typeset).}

The objective \eqref{eq:J} removes the reward for coverage by
construction: floor-level energy is pre-paid by the floor, so a fan
track that tiles the floor earns nothing and still pays curvature rent.
Whether this repair suffices in practice --- i.e.\ whether
$\Jobj(\text{fan}) > \Jobj(\text{telemetry}) > \Jobj(\text{refined})$
on the same stored trajectory sets --- is an experimental question,
answered in Sec.~\ref{sec:res-measure}
\pending{J-rescoring of the stored multi-start trajectory sets}. One
honest caveat accompanies it: we compare hypotheses by the
\emph{profiled} $\Jobj$ (nuisances minimized out, not marginalized), so
the Occam volume factor of the comb parameters is still dropped; the
floor's rent is the only coverage charge. If profiling proves
insufficient, marginalization is the principled escalation.

\subsection{The block-coordinate solver}
\label{sec:solver}

The objective \eqref{eq:J} is minimized by block-coordinate descent over
three blocks, each of which is an exactly solvable subproblem. The
alternation is a hard-EM scheme: the comb state is the latent variable,
the floor is the parameter, and the floor block is literally the M-step.

\emph{Block A (envelopes).} For fixed phases and floor, \eqref{eq:J} is
quadratic in the envelopes. The signal is demodulated per (rotor,
harmonic) by the current carrier and decimated to the envelope rate
$\fsenv$; tracks whose instantaneous frequencies approach within a
coupling radius are solved jointly (their cross terms are diagonal in
envelope time), and the resulting Hermitian positive-definite banded
system --- time-major interleaving places all couplings within a small
band --- is factorized by banded Cholesky. Whitening by the current
floor enters here: each cell's data weight is $1/S$, clamped to
$\pm15$\,dB around unity and normalized to unit geometric mean so that
the effective envelope bandwidths are floor-independent. Retaining the
cross-track coupling is what lets two tracks separated by less than a
filter bandwidth divide energy correctly instead of both claiming it ---
mandatory for interlaced quadrotor combs.

\emph{Block B (phases).} For fixed envelopes, the phase corrections are
read from the demodulated envelopes' rotation. The block measures
per-frame phase \emph{increments} of the channel-combined envelope of
each track --- increments, not absolute phases, because under the model
of Sec.~\ref{sec:model-phasenoise} increments remain informative about
the rate error long after absolute phase has decohered. Each track's
increment stream is gated by a concentration statistic (circular mean
resultant $\ge 0.5$, no near-$\pi$ steps) and weighted as $k^2 \kappa^2$
--- the Fisher weight for a rate error read at harmonic $k$, damped by
the measured concentration. The gated, weighted increments are fused and
smoothed by a Whittaker--Henderson smoother (the exact solver of the
chain-quadratic MAP subproblem; equivalently an information-form Kalman
smoother), first for the rig-common shaft correction, then per rotor,
then --- from the second outer iteration --- per strong track for
$\psi_{i,k}$.

\emph{Block C (floor).} For fixed comb state, the floor is re-estimated
from the residual: short-time spectra of $x_c$ minus the harmonic
reconstruction, with comb neighborhoods masked, smoothed by a moving
median across frequency and a cepstral lift, in a small number of time
blocks. This is the M-step; its bias (a hard-EM floor sits low by the
comb-fit posterior variance) is discussed in
Sec.~\ref{sec:discussion}.

\emph{Annealing.} The blocks alternate under a trust ladder in harmonic
order: the first iteration trusts only the coherent low harmonics
($k \le 3$) for phase reading, subsequent iterations extend to
$k \le 12$ and then to the full comb, as the improving carrier brings
higher harmonics inside the capture range of Block B
(Sec.~\ref{sec:capture} gives the arithmetic). Long recordings are
processed in overlapping windows (12\,s, 9\,s hop) and stitched with a
raised-cosine crossfade in the envelope domain.
\pending{final iteration/ladder constants after the hyperparameter
minimization pass; the current implementation uses 3 iterations with
$k$-trust $(3, 12, 80)$.}

\subsection{One machine, three readouts}
\label{sec:readouts}

The converged solver state answers all three research questions:
\begin{enumerate}
\item \emph{Decomposition}: the harmonic reconstruction (envelopes
  applied to corrected carriers) and the residual with its fitted floor
  are the two components. The validation criterion is stringent and
  visual as well as numeric: the residual must contain no evident
  harmonic structure, quantified by the \emph{excess retained} metric of
  Sec.~\ref{sec:metrics}.
\item \emph{Refinement}: the corrected rates
  $\hat r_i + \delta_i$ (with $\delta_i$ read from $D_1\theta_i$) are an
  improved trajectory estimate. Block B run under the annealing ladder
  \emph{is} the refinement stage; it is an information-form descendant
  of the phase-based VK frequency update, repaired by the increments
  formulation, the concentration gating, and the capture-matched
  ladder.
\item \emph{Measure}: the converged objective $\Jobj^\ast(\hat r)$ is a
  scalar fitness of the input trajectory set, comparable across
  hypotheses on the same window. It is the objective that the search of
  Sec.~\ref{sec:search} optimizes.
\end{enumerate}

\subsection{Hyperparameters}
\label{sec:constants}

\wip{The hyperparameter set is being actively reduced; the table below
freezes the current set and the reduction target. The design rule: every
remaining hyperparameter must be a bandwidth in hertz calibrated by the
single relation of Sec.~\ref{sec:map-objective}, a physically
interpretable threshold, or a window/rate constant.}

\begin{table}[t]
  \centering
  \caption{Frozen design constants. \pending{final minimal set after the
      principled-reduction pass}}
  \label{tab:constants}
  \scriptsize
  \setlength{\tabcolsep}{4pt}
  \begin{tabular}{@{}ll@{}}
    \toprule
    Constant & Value \\
    \midrule
    Sample / envelope rate & 16\,kHz / 62.5\,Hz \\
    Harmonics $K$; analysis window / hop & 80; 12\,s / 9\,s \\
    Shaft phase bandwidth $b_\theta$ & 1.5\,Hz \\
    Per-track bandwidth $b_\psi(k)$ & \pending{single-rule form} \\
    Envelope bandwidths $b_A(k)$ & \pending{single-rule form} \\
    Concentration gate & $\kappa \ge 0.5$ \\
    Whitening clamp & $\pm 15$\,dB \\
    Trust ladder & $k \le (3, 12, 80)$ \\
    Floor: STFT / median span / cepstral order &
      4096 / 9 bins / 40 \\
    \bottomrule
  \end{tabular}
\end{table}

\section{Blind Search}
\label{sec:search}

The solver of Sec.~\ref{sec:solver} is local: Block B reads phases of
envelopes demodulated by the \emph{current} carrier, so its capture
range is bounded (Sec.~\ref{sec:capture}), and the measure
$\Jobj^\ast$ is multimodal in the trajectories at every scale (octave
doubles, twin swaps, rational-ratio aliases). Blind operation therefore
requires a global stage whose job is mode selection: hand the local
machine an initialization inside the correct basin. This section
describes that stage; it is a magnitude-domain discrete search and is
deliberately unglamorous --- every precise instrument in this paper
lives in the local machine, and the global stage only needs to be
\emph{unable to be trapped}.

\subsection{Comb-matched dynamic programming}
\label{sec:global}

On a whitened log-magnitude spectrogram $M(f,t)$ (running-median
whitening over frequency, so combs stand out), a candidate shaft rate
$c$ is scored by the half-tooth contrast
\begin{equation}
  S(c,t) = \frac{1}{|\mathcal{K}|}\sum_{k \in \mathcal{K}}
  \Big( M(kc,\, t) - \big[M\big((k{-}\tfrac12)c,\, t\big)\big]_+ \Big),
  \label{eq:combscore}
\end{equation}
which penalizes sub-multiple aliases without the whitening-dip artifact
of a signed contrast, and frames are chained by dynamic programming over
a discretized rate grid,
\begin{equation}
  V(c,t) = S(c,t) + \max_{c'} \big[ V(c', t{-}1) - \gamma\, |c - c'|
  \big].
  \label{eq:viterbi}
\end{equation}
The recursion returns the exact global maximizer over grid paths --- no
initialization, no basins --- so its capture range is the whole grid;
its resolution is capped by grid width and continuity prior at about
$1\revps$ per steady window, which is what the local machine's capture
range must (and does) accommodate.

Blind seeding, octave disambiguation, ramp handling, and per-rotor
decoupling follow the staged design of our earlier tracking work:
constant per-rotor bases from a comb-matched scan of the time-averaged
whitened spectrum with residual re-scanning and plausibility priors; a
physical octave test (if the line at base $b$ outweighs the line at
$2b$, then $b$ is itself the blade-passage comb and is halved); a
slope-tolerant full-range coarse pass with trust gates for ramps and
takeoff; and a residual-corridor coordinate descent that decouples
individual rotors from the union template, guarded by truth-free
acceptance gates. \wip{This subsection currently summarizes; the full
description with constants is retained from the previous draft and will
be compressed into an appendix or kept inline depending on length
budget.}

\subsection{Capture range: why the machine is built this way}
\label{sec:capture}

The architecture --- global mode selection, then local phase-domain
refinement under a $k$-trust ladder, with no classical VK frequency loop
anywhere --- is dictated by one piece of arithmetic. An envelope solved
at bandwidth $b$ around an assumed tooth position only reconstructs
content within $b$ of that position; a rate error $\delta$ displaces
harmonic $k$ by $k\delta$\,Hz; hence a phase reading at harmonic $k$ is
informative only while
\begin{equation}
  k\,|\delta| \lesssim b/2,
  \label{eq:basin}
\end{equation}
and the capture radius of any fixed-band phase-domain refiner shrinks as
$1/k$ on exactly the harmonics whose Fisher weight $k^2$ favors them.
At the errors a blind global stage delivers ($0.5$--$1.5\revps$), a
fixed narrow band leaves \emph{zero} informative harmonics; the
classical VK frequency loop is inert there --- and, being an in-band
estimator, it does not fail loudly; it reports convergence. We verified
this on instrumented runs: the fused rate-error slope entering the loop's
smoother is two to three orders of magnitude below the true error, an
ablation with the loop disabled reproduces the full pipeline to
$\pm0.003\revps$, and the label-free residual is unchanged by five loop
iterations while a single increments pass improves it. The published
adaptive variant IAVKF~\cite{jiang2024iavkf} manages the same radius
implicitly by its wide-to-narrow bandwidth adaptation, with no
dependence on harmonic order; the ladder of Sec.~\ref{sec:solver} is
that schedule made explicit and $k$-aware: trust low harmonics first
(wide capture, low precision), extend $k$ as the error shrinks
(precision without loss of lock). The same arithmetic sets the division
of labor with the global stage: the ladder's entry capture radius at
$k \le 3$ accommodates the $\approx 1\revps$ resolution the DP hands
over.

\section{Experimental Setup}
\label{sec:setup}

\subsection{Data}
\label{sec:data}

\textbf{DREGON}~\cite{strauss2018dregon}: indoor free flight, 8-mic
array, resampled $44.1 \to 16$\,kHz, 1\,kHz motor telemetry; cruise twin
pairs at $0.43$ and $0.81\revps$. The dataset provides both commanded
and measured motor-speed tracks; we use the measured track
(\texttt{motors\_measured}), which is an independent sensor and thus a
legitimate external reference for blind operation, while carrying its own
audio-referenced systematic bias, which we quantify (a scale offset of
$0.35$--$0.85\%$, i.e.\ $0.3$--$0.6\revps$ at cruise, with no acoustic
lock) and quote as a floor beside every telemetry-referenced number on
this data. \textbf{Matrice~100} (recordings FLY124/FLY125): outdoor
free flight, 8-mic ring, flight-controller telemetry at
${\approx}29$\,Hz. The outdoor telemetry required auditing: a label-free
scan finds a linearly drifting clock ($R^2 = 0.94/0.92$; residual RMS
$2.9/4.5$\,ms against $12/16$\,ms for the best constant lag) and a
$0.70\%/0.71\%$ rate scale, two independent fits agreeing to $0.008$
percentage points. Corrected labels are used throughout; the correction
alone improves pooled cruise tracking error by $8.5\%$ --- label
quality, not estimator credit. The lesson we draw from both audits, and
build into the evaluation design: \emph{telemetry is a measurement, not
ground truth}, and every comparison in this paper is constructed to
survive that fact.

\subsection{Synthetic tiers}
Three synthesis tiers isolate error sources; each comprises 6 windows of
16\,s with frozen seeds. \textbf{S1}: one rotor, $K = 30$ harmonics with
the DREGON-measured amplitude rolloff, locked phases, broadband noise at
$+10$\,dB --- isolates the estimator floor. \textbf{S2}: four rotors
including one twin pair at $0.5\revps$ separation, locked phases, 0\,dB
--- isolates inter-rotor interference. \textbf{S3}: four rotors with
per-harmonic phase diffusion at the measured law and broadband noise at
$-5$\,dB --- decoherence robustness. Synthetic data is also where the
decomposition has exact ground truth: on S3 the true harmonic/broadband
split is known by construction, so the decomposition error itself
(component-wise SNR) is measurable there, not only the residual-purity
proxy.

\subsection{Protocols}
\label{sec:protocols}

\emph{Controlled imprecision.} Refiner capture ranges are compared from
identical initializations of controlled quality:
$r^{\mathrm{init}}_i(t) = r^{\mathrm{true}}_i(t) + \varepsilon (a_i +
w_i(t))$ with $a_i \sim \mathcal{N}(0,1)$ a per-rotor constant, $w_i$ a
unit-variance Ornstein--Uhlenbeck process with 1\,s correlation time,
both frozen per seed, and the single knob $\varepsilon \in \{0.1, 0.3,
0.5, 1.0, 2.0\}\revps$ sweeping from near-oracle to beyond every
refiner's capture range. All refiners, ours and baselines, receive
identical draws. Blind operation is reported both from a shared blind
initialization (isolating refinement quality) and end-to-end with each
method's native front end (the headline).

\emph{Measure validation.} The trajectory sets from the least-squares
failure experiment of Sec.~\ref{sec:whynotls} --- telemetry, refined,
blind-chain, and the coverage-fan solutions --- are frozen. Each
candidate measure (profiled least squares; $\Jobj^\ast$) scores all sets
on all windows; a measure passes if it ranks refined $\le$ telemetry $<$
fan on every window. This is a falsifiable, pre-registered test with the
failure cases chosen adversarially by an optimizer.

\subsection{Baselines}
\textbf{IAVKF}~\cite{jiang2024iavkf}: faithful reimplementation, made
multi-rotor-fair in exactly one respect (components are peeled per
(rotor, harmonic) using our track table rather than free spectral
peaks). \textbf{VKF+SWT}~\cite{li2020vkswt}: VK envelope extraction with
external synchrosqueezed-wavelet ridge IF. \textbf{Tacholess order
tracking}~\cite{peeters2019tacholess}: generalized-demodulation iterated
order spectra~\cite{olhede2005generalized}. \textbf{VK frequency loop}:
the inert classical loop, kept as an ablation row. \textbf{Profiled
least squares}: the projection fitness of Sec.~\ref{sec:whynotls}, kept
as the measure-ablation row. For the decomposition itself, the
harmonic-plus-noise baseline is frame-wise least-squares projection at
matched bandwidths \pending{decide: also HPSS-style median-filtering
separation as a model-free baseline}. Rotor-conditioned
enhancement~\cite{gulli2025rotorcond} is the \emph{consumer} of our
output, not a competitor.

\subsection{Metrics}
\label{sec:metrics}

\emph{Trajectories}: permutation-invariant mean absolute error (PIT-MAE,
\revps) against the reference labels, pooled and per rotor; capture
curves (final MAE versus $\varepsilon$); fraction of windows converged;
runtime per 16\,s window (single CPU core).

\emph{Decomposition}: the headline metric is \emph{excess retained} ---
the fraction of comb-locked energy remaining in the residual, measured
per harmonic band by an order-domain instrument: resample the residual
spectrogram to order coordinates along the estimated trajectories,
detrend the order profile by a running median (so that broadband slope
cannot masquerade as comb energy), and integrate the on-order excess
over off-order level, in bands of harmonic order ($k \le 9$, $10$--$24$,
$25$--$49$, $50$--$80$). The target is residual excess below
\pending{finalize: $\le 10\%$ per band} of the original comb energy in
every band, on every recording with telemetry. We deliberately also keep
a non-numeric criterion: the residual spectrogram must pass visual
inspection with no evident harmonic striping --- the numeric instrument
was blind to at least one failure mode (broadened combs) during
development, and we do not fully trust any single scalar here.

\emph{Downstream}: a neural noise generator conditioned on rotor speeds
is trained on labels of varying quality (telemetry; refined;
\pending{blind pseudo-labels}); the measure is whether the generator
places its learned harmonic lines on the true comb, and its validation
loss against real recordings. This is the consumer-side definition of
``sufficiently precise labels''.

\section{Results}
\label{sec:results}

\wip{This section is the working skeleton: every subsection states its
experiment, its table or figure, and --- where the measurement has
already been made in the development campaigns --- the measured anchor
numbers. Cells marked $\ast$ await the consolidated runs.}

\subsection{Decomposition on recordings with telemetry}
\label{sec:res-decomp}

\pending{The headline table: excess retained per harmonic band
($k \le 9$ / $10$--$24$ / $25$--$49$ / $50$--$80$), per recording
(DREGON, FLY124, FLY125), at telemetry and at refined labels; plus the
component-wise SNR on synthetic S3 where exact ground truth exists.
Current development state, for calibration of expectations: the low
band is essentially clean (comb depth on DREGON reduced from 3.42 to
0.31\,dB), the high bands are not yet at target. The paper claims land
only when every band passes on every recording.}

\begin{figure}[t]
  \centering
  \figplaceholder{7em}{F-D1: residual spectrograms before/after, one
    window per recording, with the excess-retained numbers inset}
  \caption{Decomposition at close-to-truth labels: original, harmonic
    reconstruction, and broadband residual. \pending{regenerate at the
    final solver configuration}}
  \label{fig:decomp}
\end{figure}

\subsection{Validation of the measure}
\label{sec:res-measure}

Two results validate $\Jobj^\ast$ as a trajectory fitness; the first is
measured, the second pending.

\emph{The least-squares failure is real.} On the three DREGON cruise
windows of the frozen protocol, unconstrained multi-start optimization
of the profiled least-squares fitness finds trajectory sets that win
that objective by $0.05$--$0.09$ while sitting $6.5$--$7.2\revps$
PIT-MAE from telemetry (fan configurations, Sec.~\ref{sec:whynotls});
on the two outdoor windows the same objective ranks telemetry first
(margins $0.19$/$0.18$) and the failure does not occur. The failure is
therefore specific to weak decohered combs --- exactly the regime the
floor model addresses.

\emph{Does $\Jobj$ fix the ranking?} \pending{J-rescoring table:
$\Jobj^\ast$ of telemetry / refined / blind chain / fan solutions on the
five frozen windows; pass criterion refined $\le$ telemetry $<$ fan
everywhere. This experiment gates every search claim downstream.}

\emph{Consistency with independent instruments.} Where the measure is
already trustworthy, it agrees with independent estimates: on six
DREGON steady windows the optimum of the fitness sits at a scale offset
of $-0.596\%$ from telemetry (range $-0.665$ to $-0.540$), inside the
confidence interval $[-0.877, -0.533]$ of a completely independent
magnitude-domain instrument --- two instruments, one displacement, which
is also consistent with the label-bias floor of Sec.~\ref{sec:data}.

\subsection{Refinement of imprecise telemetry}
\label{sec:res-refine}

\pending{Capture-curve table/figure on S1--S3 (final PIT-MAE versus
$\varepsilon$, all methods) and real-data refinement table. Measured
anchors from development: on S1 the increments pass holds an exact
initialization to $0.055\revps$ and removes a $+1.0\revps$ offset to
$0.107\revps$; on S2 the four-rotor 0\,dB floor is $\approx 0.2\revps$,
proved by a sibling-oracle arm (oracle $0.201$ vs blind $0.203$); the
iteration-harm measurement (error grows $0.167 \to 0.268$ over five
naive repeats) motivated the annealed one-machine design and will be
re-measured under the joint solver.}

\emph{The downstream test.} The consumer-side result is measured: a
noise generator trained with refined labels places its learned harmonic
lines on the true comb in the mid-order bands where the same
architecture trained with raw telemetry labels washes them out; a
constant label scale error of $0.54\%$ alone costs $8.6$\,dB of
high-order comb fidelity. Refinement is not cosmetic: it changes what a
downstream model can learn. \pending{consolidated figure: generator
spectra under raw / refined / exact labels.}

\subsection{Blind recovery}
\label{sec:res-blind}

\pending{End-to-end blind table against baselines on the frozen
protocol. Measured anchors from development: pooled blind PIT-MAE
$0.688\revps$ on DREGON steady and $1.027\revps$ on FLY124 cruise;
against the independent tachometer track on DREGON room~1 the blind
chain scores $0.97$--$1.22\revps$ on the scoreable cruise windows ---
with the reference's own $0.3$--$0.6\revps$ bias floor quoted beside
it. Baseline columns (IAVKF, VKF+SWT, tacholess OT) pending.}

\begin{figure}[t]
  \centering
  \figplaceholder{6em}{F-B1: DREGON ramp spectrogram with tracked combs
    overlaid}
  \caption{Blind tracking through a takeoff ramp.
    \pending{regenerate from final chain}}
  \label{fig:ramp}
\end{figure}

\subsection{Runtime}
\pending{Per-stage timings of the final configuration, per 16\,s window,
single CPU core; the development configuration runs the full joint
solve in low tens of seconds per window.}

\section{Discussion and Limitations}
\label{sec:discussion}

\emph{What the model does not claim.} The floor is per-microphone with a
common driver; per-rotor broadband attribution was measured to be
unidentifiable in steady flight on our arrays (wash-borne microphones,
collinear rotor speeds, non-additive powers) and is not attempted. The
per-track phase noise is treated in the window-local limit (curvature
penalties); the full stationary-OU treatment, and the spectral shape of
the per-track wobble at high $k$ (white versus correlated), are not
resolved by our measurements and are honest open questions.

\emph{Hard-EM floor bias.} The floor block is an M-step against the
current comb fit, so the converged floor sits low by the comb-fit
posterior variance; a proper-EM correction (adding the posterior
variance of the comb reconstruction to the residual power before floor
smoothing) is the principled fix. \pending{measure whether the bias is
material at the final configuration.}

\emph{Profiling versus marginalization.} $\Jobj^\ast$ compares
hypotheses after minimizing nuisances out, so the comb's Occam volume is
not charged; the floor rent is the only anti-coverage term. The
J-rescoring test is constructed to detect whether this suffices; if it
does not, the escalation path (Laplace-approximate marginalization over
envelopes, which is tractable because Block A is Gaussian) is clear,
principled, and adds no new hyperparameters.

\emph{Label bias floors.} Both reference telemetries carry measured
systematic offsets against the acoustic optimum ($0.35$--$0.85\%$ scale
on DREGON; clock drift and scale on Matrice~100, corrected before use).
All sub-$0.5\revps$ telemetry-referenced claims in this paper are quoted
beside these floors, and the downstream generator test exists precisely
because ``distance to telemetry'' stops being the right target at this
precision --- for the consumer, the acoustically-locked trajectory is
the better label.

\emph{Open.} The measured coherence budgets bound coherent integration
for any estimator; assembling them into a Cram\'er--Rao bound for blind
multi-rotor rate tracking remains open. So does the residual harmonic
excess at high orders: the target of Sec.~\ref{sec:metrics} is a
falsifiable engineering gate, not a proof that the model family
suffices, and we will report honestly if a principled configuration
cannot reach it.

\section{Conclusion}
\label{sec:conclusion}

Decomposing multirotor ego-noise into harmonic and broadband components
and recovering rotor speed trajectories from audio are one problem, and
the honest way to solve either is to solve both: a joint model in which
phase noise lives in rate increments and the broadband floor pays
likelihood rent for the energy it absorbs yields, from one
block-coordinate solver, the decomposition, a capture-matched trajectory
refinement, and --- critically --- a measure of fit that remains
trustworthy where the classical projection residual provably inverts
its ranking. On top of this machine, a deliberately simple global search
stage extends operation to fully blind recovery, evaluated against
telemetry that is itself audited rather than trusted. The decomposition
gate --- a broadband residual visibly and measurably free of harmonic
structure on every recording with telemetry --- is the standard the
system is held to, and the standard we recommend for any future claim of
drone ego-noise separation.

\bibliographystyle{IEEEtran}
\bibliography{references}

\end{document}
